\documentclass{kpi-lab}
\usepackage{tabularx}
\usepackage{biblatex}
\addbibresource{sources.bib}

\title{Лабораторна робота \textnumero 1\\
  Вибір та реалізація базових фреймворків та бібліотек
}
\author{ФБ-51мн Подолянко Тимофій}
\date{}

\begin{document}
\maketitle
\tableofcontents
\bigskip

\textbf{Мета роботи} ---
дослідити сучасні програмні бібліотеки реалізації
криптографічних механізмів,
критерії їх вибору для різних програмно-апаратних платформ
та отримати практичні навички коректного використання
криптографічних API.

\chapter{Дослідження принципів реалізації криптографічних механізмів}

При розробці прикладного програмного забезпечення не
рекомендується самостійно реалізовувати стандартні
криптографічні алгоритми, оскільки їх коректна реалізація
є складною задачею.
Криптографічні засоби захищають найбільш цінні ресурси,
тому ризик збитків через помилку значно зростає
порівняно з іншими програмними компонентами.
Поширені криптографічні бібліотеки засоби проходять аудит, процедури
детального аналізу, перевірки на стійкість до відомих класів атак.
Встановлення значень параметрів алгоритмів,
використання криптографічних примітивів
без детального їх розуміння і відповідних перевірок
створює ризик
зниження стійкості криптографічних механізмів, втрати даних.
Програмним реалізаціям також необхідно вживати заходів
задля упередження атак витоку інформації сторонніми каналами.
До того ж, використання стандартних криптографічних
бібліотек гарантує сумісність із стороннім програмним забезпеченням.

Які основні групи криптографічних примітивів та допоміжних
функцій зазвичай надають сучасні криптографічні бібліотеки?

Зазвичай сучасні криптографічні бібліотеки надають:
\begin{enumerate}
  \item Засоби генерації (криптографічно стійких)
    випадкових даних, генерації ключів.
  \item Примітиви симетричної, асиметричної криптографічної.
  \item Примітиви для реалізації односторонніх перетворень
  \item Реалізації симетричних схем шифрування, блокового і
    потокового шифрування.
  \item Реалізації асиметричних схем шифрування.
  \item Реалізації AEAD схем.
  \item Реалізації схем цифрового підпису на асиметричних примітивах.
  \item Реалізації схем кодів автентифікації повідомлень.
  \item Реалізації схем гешування, функцій одержання ключів
    (key derivation functions).
  \item Засоби серіалізації та десеріалізації даних
  \item У бібліотеках для низькорівневих мов програмування
    --- допоміжні функції для роботи з пам'яттю (zeroing,
    locking, \dots).
\end{enumerate}

Криптографічний примітив ---
низькорівневий криптографічний алгоритм, який
використовують як базовий
блок для побудови більш високорівневих криптографічних
алгоритмів \cite[NIST][]{Barker_2020}.
Протокол ---
множина правил, яка описує порядок повідомлень та структури
даних для обміну інформацією
між двома або більше сутностями, що комунікують
\cite[NIST][]{Barker_2020}.
Відповідно, криптографічний протокол використовує
криптографічні примітиви для побудови повідомлень.
Наприклад, варіанти протокола Діффі—Геллмана для узгодження ключів
потребують такі криптографічні примітиви, як
генератор випадкових чисел,
генератор простих чисел,
односторонні функції,
операції над алгебраїчними групами \cite[IETF RFC][]{rfc9846}.

Криптографічно стійкий генератор псевдовипадкових чисел ---
це алгоритм, який
дозволяє з початкових даних із достатнім рівнем ентропії
отримати велику кількість криптографічно стійких випадкових
значень \cite[RFC][]{rfc4086}.
Класичні генератори псевдовипадкових чисел спрямовані лише
на наближення статистичних властивості згенерованих послідовностей до
властивостей істинно випадкових послідовностей.
Але в криптографічно стійких послідовностях додатково
вимагається, щоб знання деяких значень з послідовності не
надавало можливості дізнатися значення інших \cite[RFC][]{rfc4086}.
\enquote{Звичайні} PRNG часто цю вимогу не задовольняють.
Таким чином, ключ, згенерований нестійким PRNG, можна було б
відновити з часткового значення.
Або з'явилася б можливість підібрати інший випадковий параметр, який
не повинен бути передбачуваним.

Низькорівневі криптографічні API зазвичай надають
можливість безпосередньо використовувати криптографічні
примітиви. Такі API самі по собі на вимагають від
розробника дотримання певного протоколу чи безпечних
практик використання даних примітивів. Неправильне
застосування низькорівневого API загрожує конфіденційності
і цілісності даних. Натомість високорівневі API реалізовані
так, щоби запобігти некоректному використанню примітивів.
Вони не повинні вимагати від користувача чи розробника
знання особливостей використаних примітивів та протоколів,
значень параметрів для безпечного застосування.

Constant-time implementation криптографічної операції
гарантує, що час виконання операції не залежить від вхідних даних.
Вона захищає від атак витоку інформації часовим каналом.

Стандарт із серії 140 Federal Information Processing
Standards FIPS 140-3 описує вимоги безпеки до проектування,
реалізації, та використання
\enquote{криптографічних модулів}, які використовуються
федеральними департаментами та відомствами (США) або їх підрядниками.
Зокрема, вимоги покривають
специфікації;
інтерфейси;
ролі, сервіси, автентифікацію;
безпеку програмного забезпечення;
операційні середовища;
фізичну безпеку;
неінвазивну безпеку;
керування чутливими параметрами безпеки;
самотестування;
гарантії життєвого циклу;
запобігання іншим атакам.
Стандарт застосовний до апаратних, програмних, комбінованих
реалізація криптографічних модулів.
Аналогом FIPS є поєднання стандартів ISO/IEC 19790б, 24759.
Статус сертифікації модуля за FIPS 140 можна
переглянути на онлайн-ресурсі програми
\href{https://csrc.nist.gov/projects/cryptographic-module-validation-program}
{CMVP}.
Наявність даної сертифікації вимагається регуляціями,
зокрема, FISMA, HIPAA, PCI DSS v4.0 та інші
(\href{https://www.encryptionconsulting.com/fips-compliance}
{дані encryptionconsulting.com}).

Стійкість RSA забезпечується складністю задачі факторизації
та бітовою довжиною чисел.
Наразі широко визнана достатньою довжина модуля в 2048 бітів,
що значно більше за розмір реєстрів звичайних
процесорів, а тому вимагає програмної реалізації
арифметичних операцій.
Окрім коректності, швидкодії, від реалізації
багаторозрядної арифметики
може вимагатися виконання операцій з постійним часом.

\chapter{Порівняння криптографічних бібліотек}

\begin{table}
  \caption{Порівняльна таблиця криптографічних бібліотек}
% tex-fmt: off
  \begin{tabularx}\textwidth{|X|p{\dimexpr \textwidth/3}|p{\dimexpr \textwidth/3}|}
    \hline
                                    & \href{https://www.pycryptodome.org/}{\textbf{pycryptodome}}                   & \href{https://libsodium.gitbook.io/doc}{\textbf{libsodium}}                   \\\hline
    Мови, платформи                 & Python (w/ C internals)                                                       & C; \emph{bindings for} Python, JVM, \dots                                     \\\hline
    Симетричне \newline шифрування  & AES, (3)DES, CAST-128, RC2, Salsa20, (X)ChaCha20, RC4; \emph{w/ modes}        & (X)Salsa20, (X)ChaCha20, AES-GCM, AEGIS                                       \\\hline
    AEAD                            & (X)ChaCha20-Poly1305; \emph{modes} GCM, CCM, EAX, OCB                         & AEGIS-(128L, 256),\newline AES256-GCM,\newline (X)ChaCha20-Poly1305(-IETF)    \\\hline
    Асим. крипто.                   & +                                                                             & +                                                                             \\\hline
    ЦП                              & RSASSA-PKCS1-v1\_5, RSASSA-PSS, (EC)DSA (FIPS 186-3,RFC6979), EdDSA           & Ed25519, Ed25519ph                                                            \\\hline
    Геш-функції                     & SHA-(1,2,3), BLAKE2(b,s), (c, Turbo)SHAKE, TupleHash, RIPE-MD160, MD5, \dots  & BLAKE2b, SipHash-2-4, SHAKE, TurboSHAKE                                       \\\hline
    KDF                             & PBKDF2, scrypt, HKDF, PBKDF1                                                  & HKDF                                                                          \\\hline
    Ініціалізація CSPRNG            & -                                                                             & +- (можна задати власну імпл.)                                                \\\hline
    Генератори випадковості         & + (спеціальний API)                                                           & (спеціальний API)                                                             \\\hline
    Високорівневий API              & -                                                                             & +                                                                             \\\hline
    Апаратне прискорення            & + (AES)                                                                       & + (AES-GCM)                                                                   \\\hline
    FIPS                            & -                                                                             & -                                                                             \\\hline
    Ліцензія                        & public domain + BSD 2-clause                                                  & \href{https://en.wikipedia.org/wiki/ISC_license}{ISC}                         \\\hline
    Вразливості                     & CVE-2018-15560 (7.5),\newline CVE-2023-52323 (5.9)                            & CVE-2025-69277 (4.5)                                                          \\\hline
    Підтримка                       & oss contributions                                                             & oss contributions                                                             \\\hline
    Активна розробка                & +                                                                             & +                                                                             \\\hline
  \end{tabularx}
% tex-fmt: on
\end{table}

\section*{Обґрунтування вибору}

Для подальшого розгляду обрано \texttt{libsodium}.
По-перше, libsodium є багатоплатформною та має біндінги,
в тому числі
\href{https://libsodium.gitbook.io/doc/bindings_for_other_languages}
{офіційні},
для багатьох поширених мов програмування.
По-друге, \texttt{libsodium} надає високорівневий API для
симетричного (потокового) шифрування, AEAD,
конфіденційного обміну повідомленнями на основі
криптографії публічних ключів,
цифрового підпису. \texttt{PyCryptodome} натомість надає
зручний API для роботи з криптографічними примітивами, що
зручно для досліджень чи експериментів, але не бажано для
звичайного практичного застосування.

\chapter{Практичне використання криптографічної бібліотеки
\texttt{libsodium}}

Для демонстрації практичного використання \texttt{libsodium}
наведемо приклади на Python із \texttt{pysodium}.

\section{Симетричне автентифіковане шифрування}

Розглянемо AEAD XChaCha20-Poly1305, реалізоване в \texttt{libsodium}.

\begin{minted}{python}
message = "тестове повідомлення".encode()
associated_data = "some very important not confidential data".encode()
secret_key = pysodium.randombytes(
    pysodium.crypto_aead_xchacha20poly1305_ietf_KEYBYTES
)
nonce = pysodium.randombytes(pysodium.crypto_aead_xchacha20poly1305_ietf_NONCEBYTES)
print(f"Message:\n\t{message.hex()} ({len(message)})")
print(f"Associated data:\n\t{associated_data} ({len(message)})")
print(f"Nonce:\n\t{nonce.hex()} ({len(nonce)})")
print(f"Key:\n\t{secret_key.hex()}")

ciphertext = pysodium.crypto_aead_xchacha20poly1305_ietf_encrypt(
    message, associated_data, nonce, secret_key
)
print(f"CT:\n\t{ciphertext.hex()} ({len(ciphertext)})")

plaintext = pysodium.crypto_aead_xchacha20poly1305_ietf_decrypt(
    ciphertext, associated_data, nonce, secret_key
)
print(f"PT:\n\t{plaintext.hex()} ({len(plaintext)})")

assert plaintext == message, "Decrypted message must be the same as the original"

modified_ct = bytearray(ciphertext)
modified_ct[len(modified_ct) // 2] = 0
modified_ct = bytes(modified_ct)
try:
    plaintext = pysodium.crypto_aead_xchacha20poly1305_ietf_decrypt(
        modified_ct, associated_data, nonce, secret_key
    )
    assert False, "Decryption of modified ciphertext must fail"
except Exception as err:
    print("Modified ciphertext could not be decrypted:", repr(err))

modified_ad = bytearray(associated_data)
modified_ad[9] = ord("!")
modified_ad = bytes(modified_ad)
try:
    print("Modified associated data:", modified_ad)
    plaintext = pysodium.crypto_aead_xchacha20poly1305_ietf_decrypt(
        ciphertext, modified_ad, nonce, secret_key
    )
    assert False, "Decryption with modified associated data must fail"
except Exception as err:
    print(
        "Won't decrypt message when associated data has been modified:", repr(err)
    )
\end{minted}
\begin{minted}{output}
Message:
        d182d0b5d181d182d0bed0b2d0b520d0bfd0bed0b2d196d0b4d0bed0bcd0bbd0b5d0bdd0bdd18f (39)
Associated data:
        b'some very important not confidential data' (39)
Nonce:
        1a66ff10ca5f6c3bc45a4b533caf0989bb6e9fc10f003c38 (24)
Key:
        6caf1fbf6384b8794fdf82573630d47ae4f2e1e96a2afc5babed44e4c96e6d7a
CT:
        50efe696d2d5125b02a1c1f6717ff256b217014a7e0753d29a8f8ae9a7b73e7d4ecd06534535ae1cdf9c16c3d47eda4c37b7d29319bece (55)
PT:
        d182d0b5d181d182d0bed0b2d0b520d0bfd0bed0b2d196d0b4d0bed0bcd0bbd0b5d0bdd0bdd18f (39)
Modified ciphertext could not be decrypted: ValueError()
Modified associated data: b'some very!important not confidential data'
Won't decrypt message when associated data has been modified: ValueError()
\end{minted}

XChaCha20-Poly1305 реалізовує AEAD схему на основі потокового шифру.
Секретний ключ використовується для шифрування (генерації гами).
Nonce виступає вектором ініціалізації для ChaCha20.
Nonce не є секретним, але повинен застосовуватися лише один раз
для збереження гарантій Poly1305.
Секретний ключ та nonce в XChaCha20-Poly1305 використовуються для
створення одноразового ключа Poly1305.
Отриманий Poly1305 authentication tag для повідомлення підтверджує
автентичність повідомлення. Цей тег автентифікує зміст
асоційованих даних та їх довжину,
шифротекст та його довжину \cite[RFC][]{rfc8439}.

В XChaCha20 Nonce має довжину 192 біти, а тому може
безпечно обиратися випадковим чином.
В ChaCha20 Nonce має лише 96 бітів,
що створює значний ризик колізіі у разі випадкової генерації,
тому для цього алгоритму рекомендовано послідовно збільшувати Nonce.
Використання пари (секретний ключ, Nonce) більше одного разу

Повторне використання Nonce для різних повідомлень
має катастрофічні наслідки для стійкості цієї схеми.
Супротивник, який спостерігає два різні повідомлення,
зашифровані з тою самою парою (секретний ключ, Nonce),
може дізнатися значення \(PT_1 \oplus PT_2 = CT_1 \oplus CT_2\).
Він також зможе створювати повідомлення з підробленим
authentication tag на даному секретному ключі.

\section{Цифровий підпис}
\begin{minted}{python}

message = "test message".encode()
public_key, secret_key = pysodium.crypto_sign_keypair()
signature = pysodium.crypto_sign_detached(message, secret_key)
print(f"""Message:
    {message}
    {message.hex()} ({len(message)})
Secret key: {secret_key.hex()} ({len(secret_key)})
Public key: {public_key.hex()} ({len(public_key)})
Signature: {signature.hex()} ({len(signature)})""")

pysodium.crypto_sign_verify_detached(
    signature,
    message,
    public_key,
)
print("Authentic message signature verification succeeds")

modified_message = bytearray(message)
modified_message[4] = ord("_")
modified_message = bytes(modified_message)
print(f"""Modified message:
    {modified_message}
    {modified_message.hex()} ({len(modified_message)})""")
try:
    pysodium.crypto_sign_verify_detached(
        signature,
        modified_message,
        public_key,
    )
    assert False, "Modified message signature verification must fail"
except ValueError as err:
    print("Modified message does not pass signature verification:", repr(err))
\end{minted}
\begin{minted}{output}
Message:
    b'test message'
    74657374206d657373616765 (12)
Secret key: 9239261ec2bced7907ea2d8df453307ed7036e4bae57e2c7729f944798bc9d5e836125631362112f05e297cfcc1041212a4aff01012604c4a5e8bde93bf8ba77 (64)
Public key: 836125631362112f05e297cfcc1041212a4aff01012604c4a5e8bde93bf8ba77 (32)
Signature: d9ecc2b007db966b940d958f975399ee77436963c8ae9c4b0e17509cebdaf5658072e78f8779af30fd914dd5244736fec5e6bd9c5ac97305dbdb8ab45a12f404 (64)
Authentic message signature verification succeeds
Modified message:
    b'test_message'
    746573745f6d657373616765 (12)
Modified message does not pass signature verification: ValueError()
\end{minted}

\chapter{Вибір криптографічної бібліотеки для різних IT-систем}

Розглянемо різні сценарії застосування криптографічних бібліотек
та вимоги і чинники, які можуть значно вплинути на вибір
бібліотеки для застосування у розробці програмного
забезпечення у кожному зі сценаріїв, та яким варто
приділити увагу першочергово.

\section{Backend-сервіс}\label{backend-service}
\begin{quote}\small\it
  Серверний застосунок обробляє конфіденційні дані
  користувачів та виконує значну кількість криптографічних операцій.
\end{quote}

\begin{enumerate}
  \item Швидкодія.
    Бібліотека має бути реалізована з
    розрахунком на великий обсяг оброблених даних.
    Наявність апаратного прискорення криптографічних операцій.
  \item Захист від атак витоку інформації сторонніми каналами.
    Зокрема, актуально, коли серверна інфраструктура
    належить третій стороні.
  \item Відповідність стандартам безпеки згідно з
    регуляторними вимогами (наприклад, FIPS-140).
  \item Протидія бінарним вразливостям.
    Оскільки сервіс публічний, то є значно вищий ризик спроб
    активної експлуатації, в тому числі для віддаленого
    виконання коду і відмови в обслуговуванні.
  \item Стійкість до атак з використанням квантових
    обчислень (post-quantum cryptography)
    Зокрема, NIST рекомендує здійснити перехід
    на квантово стійкі алгоритми до 2035 року
    \cite{Regenscheid2024}.
\end{enumerate}

\section{Мобільний застосунок}\label{mobile-app}
\begin{quote}\small\it
  Застосунок повинен локально зберігати конфіденційні дані
  користувача у зашифрованому вигляді.
\end{quote}
\begin{enumerate}
  \item Енергоефективність, ефективність обчислень і
    використання пам'яті,
    залежно від характеру роботи з конфіденційними даними у
    застосунку.
  \item Використання засобів захисту, що надаються
    відповідною операційною системою або мобільною платформою,
    та наявні апаратні засоби захисту.
  \item Портативна, багатоплатформна програмна реалізація.
    Перевага надається рішенням, використання яких однакове на
    на різних мобільних платформах і мовах програмування.
    Також рівна підтримка різних апаратних платформ є
    актуальною для Android пристроїв, що має враховувати
    різну доступність спеціалізованого обладнання,
    наборів інструкцій.
\end{enumerate}

\section{Сервіс цифрового підпису}
\begin{quote}\small\it
  Серверний сервіс використовується для створення цифрових
  підписів документів. Закритий ключ є критично важливим
  активом системи.
\end{quote}
\begin{enumerate}
  \item Підтримка Hardware Security Module обладнання.
  \item Підтримка інтеграції з Key Management System третіх сторін.
  \item Також актуальні чинники, наведені у \ref{backend-service}.
\end{enumerate}

\section{Embedded/IoT-пристрій}
\begin{quote}\small\it
  Пристрій має обмежені обчислювальні ресурси та обсяг
  пам'яті, але повинен виконувати криптографічні операції.
\end{quote}
\begin{enumerate}
  \item Мінімальний розмір.
  \item Енергоефективна, ефективна за обчислювальним навантаженням та
    використанням пам'яті реалізація алгоритмів.
  \item Реалізація спеціальних \enquote{легких}
    (lightweight cryptography) примітивів та протоколів
  \item Модульність. Можливість обирати потрібний
    функціонал (примітиви, протоколи),
    непотрібний не включати взагалі для оптимізації розміру коду.
  \item Портативність програмної реалізації.
    Існує безліч embedded апаратних архітектур,
    часто виробники надають власний SDK з компілятором
    стандартного C.
  \item Захист від витоку інформації сторонніми каналами,
    зокрема фізичними.
\end{enumerate}

\chapter{Індивідуальне завдання}

Розглянемо вибір криптографічної бібліотеки для мобільного застосунку
більш детально.

Зважаючи на вимоги і чинники, наведені у \ref{mobile-app},
для забезпечення криптографічного захисту конфіденційних даних
на пристрої користувача доцільно розпочати із застосування
криптографічних послуг обраної мобільної платформи.
У застосуванні сторонньої криптографічної бібліотеки нема
необхідності без додаткових умов,
наприклад, взаємодії між кількома користувачами.

Для конкретики розглядатимемо платформу Apple iOS.

В iOS шифрування даних застосунків реалізоване на рівні
операційної системи.
Безпечне створення, зберігання і використання ключів шифрування
ґрунтується на застосуванні апаратних засобів.
\cite[Apple Platform Security][]{apple-platform-security}
описує чотири класи захисту даних застосунків:
\begin{enumerate}
  \item A: Повний захист --- файл доступний лише коли
    пристрій в розблокованому стані
  \item B: Повний захист, за виключенням відкритих файлів
  \item C: Повний захист до першого розблокування
  \item D: Відсутній захист --- файл завжди доступний
\end{enumerate}
Встановлення параметрів захисту можливе під час використання
API файлової системи
(наприклад
  \href{https://developer.apple.com/documentation/foundation/nsdata/write(tofile:options:)}
{\texttt{-[NSData writeToFile:options:error:]}})
та за допомогою спеціальних API
(наприклад
  \href{https://developer.apple.com/documentation/foundation/nsurl/setresourcevalue(_:forkey:)}
{\texttt{-[NSURL setResourceValue:forKey:error:]}}).
Подальший доступ до захищених файлів відбувається абсолютно
прозоро для застосунку,
за тим виключенням, що доступ може бути невдалим, якщо
відповідні умови рівня захисту не виконуються.

Ключи шифрування файлів з класами захисту A, B, C
прив'язані до користувацького фактору автентифікації
(пін-коду або пароля) та унікального ключа пристрою.

Власне шифрування використовує алгоритм AES, реалізований в
апаратному модулі \enquote{Secure Enclave AES Engine}.
В сучасних пристроях цей модуль протидіє атакам
з використанням сторонніх каналів витоку інформації, зокрема,
часовим, Static Power Analysis, Dynamic Power Analysis.
Ключі, які зберігаються в AES Engine, не доступні для
програмного забезпечення, яке виконується на основному процесорі,
а також для програмного забезпечення sepOS в Secure
Enclave Processor.

Основний фреймворк для криптографічних операцій а
платформах Apple -
\href{https://developer.apple.com/documentation/security}{Security}.
API для здійснення виключно симетричної криптографії
(наприклад,
\href{https://developer.apple.com/documentation/security/seckeygeneratesymmetric(_:_:)}{SecKeyGenerateSymmetric})
наразі має статус \enquote{deprecated}.
Натомість, розробникам пропонується використовувати такі
високорівневі API, як
\href{https://developer.apple.com/documentation/security/seckeycreateencrypteddata(_:_:_:_:)}{SecKeyCreateEncryptedData},
де симетричне шифрування виконується ефемерним ключем,
захищеним через шифруванням публічним ключем відомого отримувача.

Також, в більш сучасних версіях платформ Apple розробникам
доступний фреймворк
\href{https://developer.apple.com/documentation/cryptokit}
{Apple CryptoKit},
який надає низькорівневі API для здійснення
симетричної,
асиметричної криптографії,
гешування,
автентифікації повідомлень.

Додатково розробники застосунків можуть скористатися API
\href{https://developer.apple.com/documentation/security/keychain-services}
{Keychain Services} для роботи з невеликими обсягами конфіденційної
інформації, та API
\href{https://developer.apple.com/documentation/security/certificate-key-and-trust-services}{Certificate,
Key, and Trust Servicess}
для керування
ідентичностями,
сертифікатами,
політиками,
ключами,
довірою
в інфраструктурі відкритих ключів.

Все вищенаведене показує, що використання
Data Protection та криптографічних API
платформ Apple є найбільш оптимальним рішенням за
відсутності додаткових умов.
Використання цих API надає переваги апаратного
прискорення та апаратного захисту ключів.

Високорівневі API платформ надають
високий рівень захисту даних застосунків на пристроях Apple.
Ризик некоректного використання API є мінімальним.
Зважаючи на прив'язку до пристрою, основним ризиком варто вважати
доступність даних для користувача у разі втрати пристрою.
Варто розглянути можливість механізму безпечного резервного
копіювання.
Однак, це питання виходить за межі захисту даних у застосунку.

\chapter{Висновки}

Застосування криптографічних
механізмів є критичним для безпеки інформації.
Реалізація криптографічних бібліотек повинні враховувати
безліч аспектів, які стосуються як власне криптографічних
примітивів і протоколів,
так і умов застосування реалізованих засобів,
в тому числі програмного та апаратного середовища,
використання API бібліотеки розробниками прикладного
програмного забезпечення.

Криптографічні бібліотеки, окрім програмної коректності,
повинні бути математично коректними і не порушувати властивості
криптографічних конструкцій,
а також протидіяти витоку інформації сторонніми каналами.
Нездатність задовольнити дані вимоги створює загрози
конфіденційності і цілісності інформації.
Ризики, які виникають через некоректну реалізацію криптографічних
механізмів, високі, оскільки криптографічного захисту зазвичай
потребують найцінніші ресурси.
Стандартні криптографічні бібліотеки проходять процедури аудиту,
в тому числі з боку зацікавлених організацій та спільноти,
тому їх надійність вища за довільні in-house реалізації.

Під час вибору криптографічної бібліотеки доцільно керуватися
не лише характеристиками власне бібліотеки, а й враховувати умови її
застосування. Кожне рішення пропонує певний компроміс між
зручністю, швидкодією, ступенем захисту, іншими показниками.

\printbibliography[title={Перелік використаних джерел}]

\end{document}